\documentclass[runningheads]{llncs}
\usepackage[T1]{fontenc}
\usepackage{pslatex}
\usepackage{setspace}
\usepackage{tikz}
\usepackage{graphicx}
\usepackage{color}
\usepackage{hyperref}
\usepackage{enumitem}% http://ctan.org/pkg/enumitem
\newcommand{\MoreSpace}{\spaceskip 4pt plus 1pt minus 1pt}
\newcommand{\LittleMoreSpace}{\spaceskip 3pt plus 1pt minus 1pt}
\newcommand{\LittleLessSpace}{\spaceskip 2pt plus 1pt minus 1pt}
\newcommand{\LessSpace}{\spaceskip 1.5pt plus 1pt minus 1pt}
% Default 2.5pt plus 1.5pt minus 0.6pt
\begin{document}
\title{Modern Macros: Racket's Approach to Language-Oriented Programming}
\author{Robert Bruce Findler\orcidID{0000-0002-4245-2000}}
\authorrunning{Robert Bruce Findler}
\institute{Northwestern University, Evanston IL 60208, USA\\
\email{robby@northwestern.edu}\\
\url{https://users.cs.northwestern.edu/~robby/}
}
%
\maketitle              % typeset the header of the contribution

\begin{tikzpicture}[remember picture,overlay]
    \node[xshift=125mm,yshift=-134mm,anchor=north west] at (current page.north west){%
    \includegraphics[scale=.5]{racket-logo.eps}};
\end{tikzpicture}


% new idea for the abstract, courtesy Matthias:
%   It’s the current endpoint of an evolution that started with Lisp’s
%   1964 macros.  What is it: a programmable front end. What does it
%   enable? It is used for .. notational extensions, #langs, hosted DSLs,
%   and much more. What problems does it pose: tools for
%   extensible/programmable PLs are difficult. Here are some brand-new
%   ideas.

Racket's approach to macros is the latest point in an evolution that
started in 1963 with Lisp's macros. Building on those ideas, Racket's
macros have evolved so far that, to a modern macro programmer, macros
are more helpfully understood as extending and manipulating the
compiler's front end than as a mechanism for textual substitution or
syntactic abstraction.

Having a malleable compiler front end naturally enables succicient
implementations of many languages and a look at the Racket ecosystem
reveals a wealth of examples. Scribble, a language for writing texts
uses a \LaTeX{} inspired syntax and has been used to write thousands of
pages of documentation, dozens of research papers, and at least two
books. Redex, a language for writing and testing operational
semantics, has helped numerous researchers debug their semantics and
explore their ideas. Racket's sister-language, Typed Racket, boasts
numerous type-level innovations and\\
%
{\MoreSpace{} full-fledged interoperability with untyped code. Beside these}\\
%
large efforts, Racket's macros also have enabled extensions\\
%
{\LittleMoreSpace{} on the small and medium scales as well, being the basis}\\
%
{\LessSpace{} for its pattern matcher, class system, contract system, fam-} \\
%
ily of \texttt{for} loops and more.

In this talk, I'll discuss three essential principles of\\
%
{\LessSpace{} Racket's macro system design, showing how they enable}\\
%
{\LessSpace{} language-oriented programming and provide an intellec-}\\
%
{\LessSpace{} tual foundation for understanding modern macros. More}\\
%
recently, Racket's macro system has started to move be-\\
%
yond merely the compiler's front end and I'll discuss new\\
%
{\LessSpace{} developments that build on those principles, leading towards}\\
%
nearly-automatic IDE support for programming languages in\\ 
% 
the Racket ecosystem.

\end{document}
